% TEMPLATE-MILC-v3.tex - plantilla maestra UNICA de las guias MILC v3.
%
% La usan las tres familias de guias, cada una con su driver:
%   · Grados 6-11  -> scripts/build-guias-g11.py
%   · Bebras       -> scripts/build-guias-bebras.py
%   · Semillero    -> scripts/build-guias-semillero.py
%
% Antes habia tres copias casi identicas (~400 lineas iguales, 6 distintas):
% cada mejora habia que aplicarla tres veces, y en la practica no se hacia
% (el motor de recursos vivia solo en dos de ellas). Lo que varia entre
% familias esta parametrizado en 6 placeholders que cada driver llena
% (se nombran aqui SIN su sintaxis real: el builder reemplaza por texto plano,
% tambien dentro de los comentarios, y un valor multilinea rompe el preambulo):
%
%   HEADER_IZQ        encabezado izquierdo de pagina
%   HEADER_DER        encabezado derecho de pagina
%   PORTADA_ETIQUETA  rotulo sobre el numero grande (GRADO/MOMENTO/MODULO)
%   PORTADA_SERIE     linea de serie de la portada
%   PORTADA_ID        identificador de la guia en la portada
%   PORTADA_CIRCULO   el circulito de grado (°), o vacio si no aplica
%
% El resto de claves las documenta content/guias/_SCHEMA.md. El script valida
% que no queden placeholders sin reemplazar antes de escribir el archivo final.

\documentclass[11pt,letterpaper]{article}
\usepackage[margin=2.0cm]{geometry}
\usepackage[table]{xcolor}
\usepackage{fontspec}
\usepackage[spanish]{babel}
\usepackage{tikz}
% Librerías TikZ para los diagramas de las guías (bloque `recursos:`):
%  · babel  → babel en español vuelve activos caracteres como «>», y TikZ los usa
%             en sus opciones (>=stealth, ->). Sin esta librería, cualquier
%             diagrama con flechas revienta con «\language@active@arg> has an
%             extra }». Debe ir ANTES de que se dibuje nada.
%  · positioning        → sintaxis «below left=2pt and 3pt».
%  · decorations.pathreplacing → llaves ({brace}) para señalar rangos.
\usetikzlibrary{babel,positioning,decorations.pathreplacing}
\usepackage{graphicx}
% Circuitos: el trabajo de electrónica se hace hoy con micro:bit y sus sensores
% integrados (MakeCode, sin cableado), así que no se necesitan esquemáticos.
% Si algún día entran montajes con protoboard, basta descomentar esta línea:
% los diagramas .tex/.tikz declarados en `recursos:` ya se incrustan solos.
% \usepackage{circuitikz}
\usepackage[most]{tcolorbox}
\usepackage{tabularx}
\usepackage{array}
\usepackage{fancyhdr}
\usepackage{titlesec}
\usepackage[document]{ragged2e}  % bandera a la derecha en todo el cuerpo
\usepackage{enumitem}
\usepackage{microtype}

% Helvetica Neue no trae la flecha U+2192, asi que cada «→» escrito literalmente
% en el YAML se compilaba a NADA, en silencio (462 flechas en 61 guias: rutas del
% tipo «paleta → tipografia → lockup» salian rotas). Mapeamos el caracter a su
% version matematica, que si tiene glifo. Asi el autor puede seguir escribiendo
% «→» comodamente en el YAML.
\usepackage{newunicodechar}
\newunicodechar{→}{\ensuremath{\rightarrow}}
% Mismo problema con los simbolos de marca y vinetas: el «✓» de «una palomita ✓»
% salia como cuadrito vacio en el PDF de todas las guias que lo usaban.
\usepackage{pifont}
\newunicodechar{✓}{\ding{51}}
\newunicodechar{✗}{\ding{55}}
\newunicodechar{✔}{\ding{52}}
\newunicodechar{•}{\textbullet}
\newunicodechar{≥}{\ensuremath{\geq}}
\newunicodechar{≤}{\ensuremath{\leq}}

\setmainfont{Helvetica Neue}
\setsansfont{Helvetica Neue}
\setlength{\parindent}{0pt}
\setlength{\parskip}{6pt}
% Sin particion de palabras: en 6.º-7.º una palabra cortada por guion al final
% de linea («Comuni-cacion») frena la lectura mucho mas que una linea corta.
\hyphenpenalty=10000
\exhyphenpenalty=10000
\pretolerance=2000
\tolerance=3000
\setlength{\emergencystretch}{3em}
\linespread{1.10}
\renewcommand{\arraystretch}{1.25}
\setlist{leftmargin=5mm,itemsep=1mm,topsep=1mm}
\newcolumntype{Y}{>{\RaggedRight\arraybackslash}X}

\definecolor{milcMagenta}{HTML}{E6007E}
\definecolor{milcVino}{HTML}{5A0038}
\definecolor{milcTurquesa}{HTML}{0093A5}
\definecolor{milcVerde}{HTML}{58A923}
\definecolor{milcMostaza}{HTML}{E5B400}
\definecolor{milcOcre}{HTML}{B86600}
\definecolor{milcNegro}{HTML}{241816}
\definecolor{milcGris}{HTML}{F7F7F4}
\definecolor{milcLinea}{HTML}{E8E2DD}
\definecolor{milcGradePrimary}{HTML}{FF2D87}
\definecolor{milcGradeDark}{HTML}{9F1B56}
\definecolor{milcGradeSoft}{HTML}{FFE0EE}
\definecolor{milcGradeAccent}{HTML}{CFE7FF}
\definecolor{milcGradeText}{HTML}{FFFFFF}
\color{milcNegro}

\pagestyle{fancy}
\fancyhf{}
\lhead{\small Tecnología e Informática · Grado 10}
\rhead{\small Guía 8 · MILC v3}
\cfoot{\small\thepage}
\renewcommand{\headrulewidth}{0pt}

\newtcolorbox{titlebox}[1]{
  enhanced,colback=#1,colframe=#1,arc=7mm,boxrule=0pt,
  left=7mm,right=7mm,top=3mm,bottom=3mm,
  before skip=8pt,after skip=8pt,
  fontupper=\sffamily\bfseries\Large\color{white}
}

\newtcolorbox{softbox}[3]{
  enhanced,colback=#2,colframe=#1,borderline west={4pt}{0pt}{#1},
  arc=2mm,boxrule=.4pt,left=4mm,right=4mm,top=3mm,bottom=3mm,
  before skip=6pt,after skip=8pt,title={#3},coltitle=white,
  colbacktitle=#1,fonttitle=\sffamily\bfseries,fontupper=\RaggedRight,
  attach boxed title to top left={xshift=3mm,yshift=-2mm},
  boxed title style={arc=2mm, boxrule=0pt}
}

\newtcolorbox{drawbox}[2]{
  enhanced,colback=white,colframe=#1,arc=4mm,boxrule=1pt,
  left=5mm,right=5mm,top=4mm,bottom=4mm,before skip=8pt,after skip=8pt,
  title={#2},coltitle=white,colbacktitle=#1,fonttitle=\sffamily\bfseries,
  fontupper=\RaggedRight,
  attach boxed title to top left={xshift=4mm,yshift=-2mm},
  boxed title style={arc=3mm, boxrule=0pt}
}

\newtcolorbox{infoband}[1]{
  enhanced,colback=#1!8,colframe=#1!50,arc=2mm,boxrule=.5pt,
  left=4mm,right=4mm,top=2mm,bottom=2mm,before skip=4pt,after skip=4pt,
  fontupper=\small\RaggedRight
}

% Puente narrativo entre secciones (contrato editorial Conéctate).
% Da continuidad: "Acabas de X. Ahora vas a Y porque Z."
% Visualmente: banda izquierda en color del grado, texto en itálica pequeña.
\newtcolorbox{puentebox}{
  enhanced,colback=milcGradeSoft,colframe=milcGradeDark,
  borderline west={3pt}{0pt}{milcGradeDark},
  arc=0mm,boxrule=0pt,left=5mm,right=4mm,top=2.5mm,bottom=2.5mm,
  before skip=4pt,after skip=8pt,
  fontupper=\itshape\small\color{milcNegro}\RaggedRight
}
% La flecha va en modo matematico a proposito: Helvetica Neue NO tiene el glifo
% U+2192, asi que \textrightarrow se compilaba a NADA («Missing character: There
% is no → in font Helvetica Neue») y el puente salia sin flecha. En modo
% matematico la flecha viene de la fuente matematica y si se dibuja.
\newcommand{\puente}[1]{%
  \begin{puentebox}%
    \textcolor{milcGradeDark}{$\rightarrow$}\hspace{2mm}#1%
  \end{puentebox}%
}

\newcommand{\linead}[1]{\noindent\color{milcLinea}\rule{#1}{0.5pt}\color{milcNegro}}
\newcommand{\respuesta}[1]{\par\vspace{2mm}\foreach \n in {1,...,#1}{\noindent\color{milcLinea}\rule{\linewidth}{0.5pt}\par\vspace{5mm}}\color{milcNegro}}
\newcommand{\checkbox}{\textcolor{milcGradeDark}{\fbox{\phantom{x}}}\hspace{5pt}}

\newcommand{\phasecard}[4]{
\begin{tcolorbox}[enhanced,colback=#3,colframe=#2,arc=3mm,boxrule=.5pt,height=3.05cm,valign=center,halign=center,left=1.5mm,right=1.5mm,top=2mm,bottom=2mm]
{\sffamily\bfseries\fontsize{22}{24}\selectfont\textcolor{#2}{#1}}\par
{\sffamily\bfseries\small #4}
\end{tcolorbox}
}

% ── Piezas del rediseno 2026 (legibilidad 6.º-11.º) ───────────────────
% Banda de estacion: reemplaza el titlebox macizo. Titulo a la izquierda,
% dato practico (tiempo/modalidad) a la derecha, en una sola linea.
\newcommand{\milcestacion}[2]{%
\begin{tcolorbox}[enhanced,colback=#1,colframe=#1,arc=2mm,boxrule=0pt,
  left=5mm,right=5mm,top=2.6mm,bottom=2.6mm,before skip=11pt,after skip=7pt]
{\sffamily\bfseries\fontsize{15}{18}\selectfont\color{white}#2}
\end{tcolorbox}}

% Tarjeta de paso numerada: sustituye las tablas «Pilar / Aplicacion», que
% eran formato de consulta docente, no de estudiante. El numero va en un
% circulo sobre el borde para que la secuencia se lea de un vistazo.
\newcommand{\milcpaso}[3]{%
\begin{tcolorbox}[enhanced,colback=white,colframe=milcLinea,arc=1.5mm,boxrule=.9pt,
  left=14mm,right=4mm,top=3mm,bottom=3mm,before skip=4pt,after skip=4pt,
  fontupper=\RaggedRight,
  overlay={\node[anchor=north west,circle,fill=#2,text=white,inner sep=0pt,
    minimum size=7.4mm,font=\sffamily\bfseries\small]
    at ([xshift=3.6mm,yshift=-3.2mm]frame.north west) {#1};}]
#3
\end{tcolorbox}}

% Casilla marcable de tamano util con lapiz (la anterior era \fbox{\phantom{x}},
% de alto variable segun el texto que la seguia).
\newcommand{\milcmarca}{\framebox[3.8mm]{\rule{0pt}{3.4mm}}}

% Fila marcable de lista suelta (checks de la estacion 1).
\newcommand{\milccheckrow}[1]{%
\par\vspace{1.8mm}\noindent
\makebox[6.4mm][l]{\raisebox{-.55mm}{\textcolor{milcNegro}{\milcmarca}}}%
\parbox[t]{\dimexpr\linewidth-6.4mm\relax}{\RaggedRight #1}\par}

% Celda marcable dentro de la rubrica (tabla de autoevaluacion). Misma
% sangria francesa, pero pensada para vivir dentro de una columna X.
\newcommand{\milcopcion}[1]{%
\makebox[5.2mm][l]{\raisebox{-.55mm}{\textcolor{milcNegro}{\milcmarca}}}%
\parbox[t]{\dimexpr\linewidth-5.2mm\relax}{\RaggedRight #1}}

% ── Recursos visuales de la guía ──────────────────────────────────────
% Declarados en el bloque `recursos:` del YAML y emitidos por el builder.
% \guiaFigura   → imagen rasterizada o PDF (foto, carta celeste, montaje).
% \guiaDiagrama → fuente TikZ/circuitikz (.tex) incrustada como vector:
%                 es la vía para circuitos y esquemáticos editables.
% #1 = ruta relativa al .tex · #2 = pie (puede ir vacío)
\newcommand{\guiaFigura}[2]{%
\begin{center}
\includegraphics[width=0.88\linewidth,height=0.38\textheight,keepaspectratio]{#1}\\[1.5mm]
{\footnotesize\itshape\textcolor{milcNegro!75}{#2}}
\end{center}
}

\newcommand{\guiaDiagrama}[2]{%
\begin{center}
\input{#1}\\[1.5mm]
{\footnotesize\itshape\textcolor{milcNegro!75}{#2}}
\end{center}
}

\begin{document}
\thispagestyle{empty}

% ── Portada ───────────────────────────────────────────────────────────
% Antes: un rectangulo de color a pagina completa. Se veia bien en pantalla,
% pero en la fotocopiadora del colegio se comia el toner y salia gris sucio.
% Ahora la banda ocupa el tercio superior y el resto va en blanco.
\begin{tikzpicture}[remember picture,overlay]
  \path[fill=milcGradePrimary]
    (current page.north west) rectangle ([yshift=-10.6cm]current page.north east);

  \node[anchor=north west,text=milcGradeText] at ([xshift=2.0cm,yshift=-1.55cm]current page.north west)
    {\sffamily\bfseries\fontsize{12}{14}\selectfont GRADO};
  \node[anchor=north west,text=milcGradeText,opacity=.85] at ([xshift=2.0cm,yshift=-2.35cm]current page.north west)
    {\sffamily\bfseries\fontsize{8.5}{10}\selectfont SERIE GUÍAS · TECNOLOGÍA E INFORMÁTICA};
  \node[anchor=north east,text=milcGradeText,opacity=.85,align=right] at ([xshift=-2.0cm,yshift=-1.55cm]current page.north east)
    {\sffamily\bfseries\fontsize{9}{12}\selectfont I.E. Sor María Juliana\\Cartago, Valle del Cauca};

  \node[anchor=west,text=milcGradeText] (gradonum) at ([xshift=1.85cm,yshift=-7.1cm]current page.north west)
    {\sffamily\bfseries\fontsize{112}{112}\selectfont 10};
  % El circulito de grado (°) solo aplica a las guias de grado («11°»); en
  % Bebras y Semillero el numero grande es un momento/modulo, no un grado.
  % Se posiciona relativo al nodo `gradonum`, no en coordenadas absolutas.
  \draw[milcGradeText,line width=1.6pt]
    ([xshift=2.0mm,yshift=-4.5mm]gradonum.north east) circle (.15cm);

  \node[anchor=west,text=milcGradeText,text width=11.4cm,align=left] at ([xshift=1.5cm,yshift=0.5cm]gradonum.east)
    {
     {\sffamily\bfseries\fontsize{9}{11}\selectfont\MakeUppercase{Período 1 · Escritura abierta con IA}}\\[3mm]
     {\sffamily\bfseries\fontsize{21}{25}\selectfont\setlength{\baselineskip}{27pt}Diagramación ---\\[4pt]Canva, Google Docs\\[4pt]y LibreOffice}\\[3.5mm]
     {\sffamily\bfseries\fontsize{15}{18}\selectfont Guía 8-10-TIC}
    };
\end{tikzpicture}

\vspace*{9.4cm}
\vfill

{\sffamily\bfseries\fontsize{8}{10}\selectfont\textcolor{milcGradeDark}{LO QUE TE LLEVAS HOY}}

\begin{tcolorbox}[enhanced,colback=white,colframe=milcNegro,arc=2mm,boxrule=1.6pt,
  left=5mm,right=5mm,top=4mm,bottom=4mm,before skip=3pt,after skip=12pt,
  fontupper=\RaggedRight\fontsize{11}{15.5}\selectfont]
Documento maquetado de tu libro con (a) plantilla profesional elegida y aplicada en Canva, Google Docs o LibreOffice, (b) tipografía coherente con 1-2 familias máximo (1 para títulos, 1 para cuerpo), (c) jerarquía visual clara con 4 niveles (título de libro, título de capítulo, subtítulo, cuerpo), (d) numeración de páginas automática y tabla de contenido autogenerada, (e) capítulo de muestra completamente diagramado de al menos 6 páginas con la maquetación final aplicada
\end{tcolorbox}

\begin{tcolorbox}[enhanced,colback=milcGris,colframe=milcGris,arc=2mm,boxrule=0pt,
  left=5mm,right=5mm,top=3.5mm,bottom=3.5mm,after skip=12pt]
{\sffamily\bfseries\fontsize{8}{10}\selectfont\textcolor{milcNegro!55}{ESTA GUÍA ES DE}}\\[3mm]
\begin{tabularx}{\linewidth}{X p{3.2cm} p{3.2cm}}
\linead{\linewidth} & \linead{3.0cm} & \linead{3.0cm}\\[-1mm]
{\fontsize{8.5}{10}\selectfont\textcolor{milcNegro!55}{Nombre completo}} &
{\fontsize{8.5}{10}\selectfont\textcolor{milcNegro!55}{Grupo}} &
{\fontsize{8.5}{10}\selectfont\textcolor{milcNegro!55}{Fecha}}\\
\end{tabularx}
\end{tcolorbox}

\vfill

\noindent\textcolor{milcLinea}{\rule{\linewidth}{1pt}}\\[2mm]
{\sffamily\bfseries\fontsize{9.5}{12}\selectfont Autor: Dr. Álvaro Cárdenas Orozco}\\[1mm]
{\sffamily\fontsize{8.5}{11}\selectfont\textcolor{milcNegro!55}{Metodología MILC · apertura ancestral · pensamiento computacional · triángulo Dussel–estoicismo–Floridi}}

\newpage

\section*{Apertura: Diagramación con herramientas gratuitas --- Canva, Google Docs, LibreOffice}

\begin{softbox}{milcGradeDark}{milcGradeSoft}{Saber ancestral}
En el centro de Cartago, en las imprentas pequeñas que sostuvieron las publicaciones locales antes de la era digital, había un oficio invisible para el cliente pero crucial para la calidad del libro: \textbf{el imprentero que organizaba los tipos}. Los \emph{tipos} eran las piezas físicas de metal, cada una con una letra grabada. Antes de imprimir cualquier página, el imprentero tomaba los tipos uno por uno y los acomodaba en una caja llamada \emph{platina}, en el orden exacto del texto. Si quería poner \emph{LIBRO}, tomaba L, I, B, R, O y los alineaba. Si los tipos quedaban desordenados, la página salía \emph{LBRIO} o peor. Cada letra tenía su cajón en una bandeja grande: las mayúsculas en cajones grandes (de ahí viene el nombre \emph{caja alta}), las minúsculas en cajones pequeños (\emph{caja baja}). El imprentero memorizaba la posición de cada letra: podía componer una línea sin mirar. La disciplina era estricta: \textbf{cada letra en su lugar, cada renglón en su orden, cada página alineada con la siguiente}. Si fallaba, el libro completo salía mal. Esa práctica ancestral de organización tipográfica es la misma disciplina que la diagramación moderna pide con herramientas digitales gratuitas.
\end{softbox}

\begin{softbox}{milcTurquesa}{milcGris}{Saber contemporáneo}
La \textbf{diagramación} (también llamada maquetación, layout o tipografía aplicada) es el arte de organizar texto e imagen en páginas para que sean legibles, coherentes y agradables. En 2026, hay 3 herramientas gratuitas que cubren la mayoría de necesidades: (1) \textbf{Canva} (\texttt{canva.com}): el más fácil para principiantes. Tiene plantillas de libros gratuitas, drag-and-drop, exportación a PDF. Ideal si tu libro tiene componente visual fuerte (manuales, libros ilustrados, fanzines). (2) \textbf{Google Docs} (\texttt{docs.google.com}): integrado con tu cuenta, colaborativo, exportación a PDF. Bueno para libros mayormente de texto (novelas, ensayos, memorias). Tiene estilos, tabla de contenido automática, numeración de páginas. (3) \textbf{LibreOffice Writer}: gratis, descargable, similar a Word. Bueno para control fino de la maquetación y para usuarios que prefieren trabajar offline. La regla profesional: \textbf{elige la herramienta según el tipo de libro, no según la moda}. Para tu libro de 80 páginas, las 3 funcionan; lo importante es aplicar bien los \textbf{4 principios fundamentales}: tipografía coherente, jerarquía visual clara, espacios respirables, numeración automática.
\end{softbox}

\begin{softbox}{milcMostaza}{milcGris}{Pregunta-puente}
\textit{¿Qué sabía el imprentero al acomodar cada tipo en su cajón con disciplina, que el novato olvida cuando entrega un libro con 5 tipografías mezcladas y márgenes inconsistentes? ¿Y por qué los 4 principios de diagramación valen más que las plantillas elegantes mal aplicadas?}
\end{softbox}

\begin{softbox}{milcVerde}{milcGris}{Saber y saber hacer}
Al terminar podrás: (1) \textbf{identificar} cuál herramienta gratuita conviene a tu tipo de libro; (2) \textbf{aplicar} los 4 principios fundamentales de diagramación (tipografía coherente, jerarquía visual, espacios respirables, numeración automática); (3) \textbf{crear} la maquetación completa del capítulo de muestra con todos los elementos; (4) \textbf{evaluar} la legibilidad y coherencia de tu diagramación antes de pasar a la producción final.
\end{softbox}



\small
\begin{tabularx}{\textwidth}{>{\bfseries}p{3.0cm}Y}
\rowcolor{milcGradePrimary!12}Autor & Dr. Álvaro Cárdenas Orozco\\
DBA & Producción editorial mediada por IA con criterio ético y técnico (MEN, T\&I)\\
Referentes & Ley 23 de 1982 · Floridi (ética de la IA generativa) · Dussel (autoría con dignidad)\\
Producto final & Documento maquetado de tu libro con (a) plantilla profesional elegida y aplicada en Canva, Google Docs o LibreOffice, (b) tipografía coherente con 1-2 familias máximo (1 para títulos, 1 para cuerpo), (c) jerarquía visual clara con 4 niveles (título de libro, título de capítulo, subtítulo, cuerpo), (d) numeración de páginas automática y tabla de contenido autogenerada, (e) capítulo de muestra completamente diagramado de al menos 6 páginas con la maquetación final aplicada\\
\end{tabularx}
\normalsize

\puente{Antes de empezar, mira el viaje completo. Tienes texto V3, estructura cerrada, derechos definidos, portada lista. Hoy le pones \textbf{cuerpo visual} al libro: tipografía, márgenes, jerarquía, numeración. La sesión 9 (producción final) integrará todo en el PDF final de 80 páginas.}

\milcestacion{milcGradeDark}{Ruta MILC: así vamos a aprender}
\begin{tabularx}{\textwidth}{>{\centering\arraybackslash}X>{\centering\arraybackslash}X>{\centering\arraybackslash}X>{\centering\arraybackslash}X}
\phasecard{1}{milcVerde}{milcGris}{Escucha\\[-1mm]\small Observo, escucho y pregunto.} &
\phasecard{2}{milcTurquesa}{milcGris}{Entender\\[-1mm]\small Sistematizo y ordeno.} &
\phasecard{3}{milcMagenta}{milcGris}{Hacer\\[-1mm]\small Construyo y aplico.} &
\phasecard{4}{milcVino}{milcGris}{Cierre\\[-1mm]\small Evalúo y reflexiono.}
\end{tabularx}


\puente{Antes de teorizar, vas a hacer una \emph{auditoría rápida} de un libro físico que tengas a mano (de tu casa, biblioteca o aula): anota qué tipografía usa para títulos vs cuerpo, qué tamaño tienen los márgenes, cómo numera las páginas, qué tipos de jerarquía visual hay (capítulo, subtítulo, etc.).}

\milcestacion{milcVerde}{Estación 1 de 3 · Escucha}

\begin{softbox}{milcVerde}{milcGris}{Escena para escuchar}
\textbf{Actividad 1 · IDENTIFICA --- Auditoría visual de un libro físico} (15 min · individual). Toma un libro físico que tengas a mano. Sin leer el contenido, examina la maquetación. Anota: (1) ¿qué tipografía usa para títulos de capítulo? ¿es la misma del cuerpo o distinta? (2) ¿qué tamaño aproximado tienen los márgenes (estrecho, mediano, ancho)? (3) ¿cómo numera las páginas (abajo centrado, arriba al lado, sin numeración)? (4) ¿cuántos niveles de jerarquía ves (título de libro, capítulo, subtítulo, cuerpo, pie de página)? (5) ¿hay encabezado con título del capítulo? Esa observación intuitiva te enseña el oficio de la maquetación.
\end{softbox}

{\sffamily\bfseries\fontsize{8}{10}\selectfont\textcolor{milcNegro!55}{MARCA CUANDO LO HAYAS HECHO}}
\milccheckrow{Examiné un libro físico con criterio visual.}
\milccheckrow{Identifiqué tipografía, márgenes, numeración, jerarquía.}
\milccheckrow{Reconozco que la maquetación se nota cuando funciona bien.}

\begin{infoband}{milcVerde}


\textbf{Lo que va al cuaderno (Actividad 1 · IDENTIFICA).} Título: «Actividad 1 --- Auditoría visual de un libro». \textbf{Formato:} ficha con 5 datos observados + 1 línea sobre qué hace que se sienta profesional. \textbf{Sabes que terminaste cuando} reconoces los principios invisibles que sostienen un libro bien hecho.
\end{infoband}

\puente{Ya tienes referencia visual. Ahora vas a entender la \textbf{anatomía de los 4 principios de diagramación}: tipografía coherente, jerarquía visual, espacios respirables, numeración automática. Y los 5 errores típicos del editor novato.}

\milcestacion{milcTurquesa}{Estación 2 de 3 · Entender}

\begin{softbox}{milcTurquesa}{milcGris}{Lo que vas a entender}
Tu auditoría te da vocabulario visual. Ahora necesitas la \textbf{anatomía formal} de los 4 principios de diagramación profesional.
\end{softbox}

\milcpaso{1}{milcTurquesa}{Los \textbf{4 principios fundamentales} de diagramación se descomponen así: (1) \textbf{Tipografía coherente}: máximo 2 familias en todo el libro. Una para títulos (puede ser sans-serif o serif distintiva, ej. Inter, Montserrat, Playfair), otra para cuerpo (típicamente serif legible, ej. Georgia, Lora, Merriweather). Tamaños tipo: título libro 28-36pt, capítulo 18-22pt, subtítulo 14-16pt, cuerpo 11-12pt, pie 9-10pt. (2) \textbf{Jerarquía visual clara}: cada nivel debe distinguirse a simple vista por tamaño, peso (bold) o color. El lector debe poder ubicarse sin esfuerzo. (3) \textbf{Espacios respirables}: márgenes de 2-2.5cm en libros, interlineado de 1.3-1.5, espacios entre párrafos. La página saturada cansa. (4) \textbf{Numeración automática}: nunca a mano. En Docs y LibreOffice: Insertar → Número de página. En Canva: usar plantilla que ya la tenga.}
\milcpaso{2}{milcTurquesa}{La \textbf{tabla de contenido automática} se genera en Docs/LibreOffice así: (a) marca cada título de capítulo con estilo \emph{Encabezado 1}. (b) marca subtítulos con \emph{Encabezado 2}. (c) ve al inicio del documento → Insertar → Tabla de contenidos. La tabla se genera sola y se puede actualizar si añades capítulos. En Canva la tabla de contenido se diseña a mano pero los enlaces internos también se pueden usar.}
\milcpaso{3}{milcTurquesa}{Las \textbf{plantillas profesionales} en Canva: busca \emph{``book interior''}, \emph{``ebook template''}, \emph{``novel layout''}. Hay decenas gratuitas. En Docs: usa los estilos predefinidos (Heading 1, Heading 2, Normal text) y configura márgenes en Archivo → Configuración de página. En LibreOffice: similar a Docs pero con más control fino de paginación.}
\milcpaso{4}{milcTurquesa}{Algoritmo: (1) elegir herramienta según tipo de libro. (2) abrir plantilla o documento nuevo. (3) configurar tamaño de página (típico A5 o carta personalizada). (4) elegir 2 tipografías (una para títulos, una para cuerpo). (5) aplicar estilos a los 4 niveles. (6) configurar márgenes y interlineado. (7) insertar numeración automática. (8) generar tabla de contenido. (9) diagramar capítulo de muestra completo.}

\begin{drawbox}{milcTurquesa}{Anatomía de la diagramación}
\textbf{Tipografía:} máximo 2 familias. \\[1mm] \textbf{Jerarquía:} 4 niveles distinguibles. \\[1mm] \textbf{Espacios:} márgenes + interlineado respirables. \\[1mm] \textbf{Numeración:} automática siempre. \\[1mm] \textbf{Tabla de contenido:} autogenerada con estilos.
\end{drawbox}

\begin{softbox}{milcMostaza}{milcGris}{Errores comunes y su corrección}
(1) \textbf{Mezcla de tipografías}: 4-5 fuentes distintas hacen el libro caótico. (2) \textbf{Sin jerarquía clara}: el lector no sabe si está en capítulo nuevo o sigue en el anterior. (3) \textbf{Páginas saturadas}: márgenes mínimos, interlineado apretado, texto que asfixia. (4) \textbf{Numeración manual}: si añades una página, todo el resto queda mal. (5) \textbf{Sin tabla de contenido}: el lector no puede saltar entre capítulos.
\end{softbox}

\begin{infoband}{milcTurquesa}


\textbf{Lo que va al cuaderno (Actividad 2 · APLICA).} Título: «Actividad 2 --- Anatomía diagramación». \textbf{Formato:} (a) los 4 principios con valores típicos; (b) las 3 herramientas con su uso; (c) los 5 errores con marca del que más temes. \textbf{Sabes que terminaste cuando} podrías diagramar otro libro sin abrir esta guía.
\end{infoband}

\puente{Tienes el método. Toca aplicarlo: eliges herramienta, aplicas plantilla profesional, configuras tipografía, defines jerarquía, agregas numeración, diagramas el capítulo de muestra completo.}

\milcestacion{milcMagenta}{Estación 3 de 3 · Hacer}

\begin{softbox}{milcMagenta}{milcGris}{Cómo vas a construir tu producto}
\textbf{Actividad 3 · CREA + EVALÚA --- Diagramación del capítulo de muestra} (40 min en sesión + tiempo en casa · individual con dispositivo). Hora de aplicar. Eliges herramienta, configuras plantilla, diagramas el capítulo de muestra completo de 6+ páginas con todos los elementos.
\end{softbox}

\milcpaso{1}{milcMagenta}{Tu diagramación se construye en 7 pasos: (1) elegir herramienta (recomiendo Google Docs para libros mayormente de texto, Canva para libros visuales). (2) abrir plantilla o crear documento nuevo configurando tamaño A5 o carta personalizada. (3) elegir tipografía: 1 para títulos, 1 para cuerpo. (4) aplicar estilos a los 4 niveles. (5) configurar márgenes (2-2.5cm) e interlineado (1.3-1.5). (6) insertar numeración automática. (7) diagramar el capítulo de muestra de 6+ páginas con todos los elementos visibles.}
\milcpaso{2}{milcMagenta}{Sigue el patrón \textbf{``coherencia sobre variedad''}: el lector se beneficia de la coherencia visual, no de la variedad. Si tu capítulo 1 tiene tipografía X y tu capítulo 5 tipografía Y, el lector se desconcentra. La phronesis del diagramador consiste en \textbf{aplicar las mismas reglas a todo el libro sin excepción}.}
\milcpaso{3}{milcMagenta}{El capítulo de muestra completo debe demostrar: (a) título de capítulo con estilo definido. (b) al menos 1 subtítulo intermedio. (c) cuerpo de texto con interlineado correcto. (d) numeración de páginas funcionando. (e) encabezado de página si lo usas. (f) espacios respirables al inicio y fin del capítulo. (g) ortografía y puntuación revisadas.}
\milcpaso{4}{milcMagenta}{Algoritmo: (1) elegir herramienta. (2) configurar plantilla. (3) aplicar tipografía. (4) aplicar estilos. (5) configurar márgenes. (6) insertar numeración. (7) diagramar 6+ páginas del capítulo. (8) verificar legibilidad en pantalla y en PDF exportado.}

\begin{drawbox}{milcMagenta}{Checklist antes de entregar}
\checkbox Herramienta elegida y plantilla aplicada. \\[1mm] \checkbox 1-2 tipografías máximo (no más). \\[1mm] \checkbox 4 niveles de jerarquía visible. \\[1mm] \checkbox Márgenes 2-2.5cm e interlineado 1.3-1.5. \\[1mm] \checkbox Numeración de páginas automática. \\[1mm] \checkbox Tabla de contenido autogenerada (si aplica). \\[1mm] \checkbox Capítulo de muestra de 6+ páginas diagramado. \\[1mm] \checkbox PDF exportado y revisado.
\end{drawbox}

\begin{softbox}{milcMagenta}{milcGris}{Plantilla de trabajo}
\textbf{Herramienta.} \textit{Canva / Docs / LibreOffice.} \\[1mm] \textbf{Tipografía.} \textit{1 para títulos, 1 para cuerpo.} \\[1mm] \textbf{Jerarquía.} \textit{4 niveles con estilos.} \\[1mm] \textbf{Espacios.} \textit{Márgenes + interlineado.} \\[1mm] \textbf{Numeración.} \textit{Automática.} \\[1mm] \textbf{TOC.} \textit{Autogenerada con estilos.} \\[1mm] \textbf{Capítulo muestra.} \textit{6+ páginas diagramadas.}
\end{softbox}

\begin{infoband}{milcMagenta}


\textbf{Lo que va al cuaderno (Actividad 3 · CREA + EVALÚA).} Título: «Actividad 3 --- Capítulo diagramado». \textbf{Formato:} (a) referencia al archivo Docs/Canva/LibreOffice; (b) capturas del capítulo diagramado; (c) decisiones tipográficas explicadas; (d) PDF exportado del capítulo de muestra. \textbf{Sabes que terminaste cuando} tu capítulo se ve como página de libro profesional.
\end{infoband}

\puente{Lo que sigue resume \emph{qué} entregas hoy: documento con plantilla + tipografía + jerarquía + numeración + capítulo diagramado. El criterio principal: que tu capítulo de muestra se vea como página de libro profesional, no como documento Word casual.}

\milcestacion{milcMostaza}{Producto de la sesión}
\textbf{Documento maquetado + capítulo de muestra de 6+ páginas + PDF}.
\begin{softbox}{milcMostaza}{milcGris}{Criterios de calidad}
(1) Herramienta gratuita aplicada correctamente. (2) 1-2 tipografías coherentes. (3) 4 niveles de jerarquía visible. (4) Márgenes y interlineado respirables. (5) Numeración automática funcionando. (6) Capítulo de muestra de 6+ páginas. (7) PDF exportado para revisión.
\end{softbox}

\puente{Antes de cerrar, mira la diagramación desde las cinco dimensiones humanas. Maquetar con disciplina es respeto por el lector que verá tu libro; entregar texto crudo sin diagramación es desprecio del oficio editorial.}

\milcestacion{milcVino}{Evaluación liberadora · cinco dimensiones}

\begin{softbox}{milcVino}{milcGris}{Sentido de la evaluación}
Evaluar no es solo revisar una nota. Es reconocer qué aprendí, qué cambió en mí, qué emociones manejé, cómo me relaciono con la ciudad y la comunidad, y qué le dejaré a quien viene detrás.
\end{softbox}

\textbf{Mi reflexión liberadora en cinco dimensiones.} Completa cada frase iniciada:

\small
\begin{tabularx}{\textwidth}{>{\bfseries\RaggedRight\arraybackslash}p{3.4cm}Y>{\RaggedRight\arraybackslash}p{4.6cm}}
\rowcolor{milcVino}\color{white}Dimensión & \color{white}\bfseries Mi respuesta (frame starter) & \color{white}\bfseries Algo que descubrí\\
1. Personal & Lo que más cambió en mí fue \linead{6.5cm} & \linead{4.2cm}\\
2. Emocional & La emoción más fuerte fue \linead{2.5cm} y la regulé \linead{3cm} & \linead{4.2cm}\\
3. Ciudadana & En democracia esto importa porque \linead{6cm} & \linead{4.2cm}\\
4. Local (Cartago) & En mi barrio sería útil para \linead{6.5cm} & \linead{4.2cm}\\
5. Intergeneracional & A mi abuela/abuelo le diría \linead{6.5cm} & \linead{4.2cm}\\
\end{tabularx}
\normalsize

\begin{infoband}{milcVino}
\textbf{Para la próxima sustentación.} Vuelve a esta tabla en seis meses. Verás que las respuestas cambian — no porque lo aprendido cambie, sino porque tú cambias. La evaluación liberadora es una conversación contigo a través del tiempo.
\end{infoband}


\puente{Y para terminar, tres voces sobre la organización visual. Dussel sobre los libros que respetan al lector apurado, Marco Aurelio sobre la coherencia tipográfica, Floridi sobre la diagramación como ética de la legibilidad.}

\milcestacion{milcGradeDark}{Triángulo de pensamiento · cierre filosófico}

\begin{softbox}{milcVino}{milcGris}{Dussel · lente del nosotros}
\textit{La diagramación clara respeta al lector apurado; la página saturada lo expulsa antes de leer.}

La mayoría de tus lectores leerán tu libro en condiciones imperfectas: en el bus, antes de dormir, entre clases. La filosofía de la liberación pide diseñar para esas condiciones: tipografía legible, espacios respirables, jerarquía clara. La página saturada exige condiciones ideales que el lector real no tiene. La phronesis del diagramador consiste en \textbf{respetar el contexto real de lectura, no el ideal académico}.

\textbf{¿Mi diagramación respeta al lector en condiciones reales, o asume lector ideal con tiempo ilimitado?}
\end{softbox}
\linead{\linewidth}\\[1mm]
\linead{\linewidth}

\begin{softbox}{milcMostaza}{milcGris}{Estoicismo · Marco Aurelio · cuidado interior}
\textit{La coherencia tipográfica es virtud editorial; la mezcla caprichosa es debilidad disfrazada de creatividad.}

Es tentador usar varias tipografías \emph{``para que se vea variado''}. La disciplina estoica recomienda lo opuesto: \emph{coherencia sobre variedad}. Una sola familia para títulos y otra para cuerpo bastan para todo el libro. La variedad innecesaria desconcentra al lector. La phronesis del oficio honra la consistencia.

\textbf{¿Mi diagramación es coherente en todo el libro, o mezclé tipografías por capricho?}
\end{softbox}
\linead{\linewidth}\\[1mm]
\linead{\linewidth}

\begin{softbox}{milcTurquesa}{milcGris}{Floridi · lente de la infoesfera}
\textit{La legibilidad bien diseñada es ética informacional en la era de la atención escasa.}

En la era de la información, los lectores tienen miles de opciones de qué leer. La ética editorial pide que cuando alguien decide leer tu libro, la maquetación facilite la lectura, no la complique. Tu diagramación de hoy te entrena en esa práctica que rige toda producción editorial profesional: \textbf{cuidar la legibilidad como respeto por el lector}.

\textbf{¿Mi diagramación facilita la lectura, o pide al lector trabajo extra que no debería hacer?}
\end{softbox}
\linead{\linewidth}\\[1mm]
\linead{\linewidth}

\begin{infoband}{milcNegro}
\textbf{Nota para el docente.} Las tres formulaciones del triángulo son la síntesis del autor de la guía, no citas textuales de los pensadores. Si vas a citarlos en clase, remite a la obra original.
\end{infoband}

\milcestacion{milcVino}{Mi autoevaluación · marca una casilla por fila}

{\small
\setlength{\extrarowheight}{2.2pt}
\begin{tabularx}{\textwidth}{>{\RaggedRight\arraybackslash\bfseries}p{3.0cm}YYY}
\rowcolor{milcVino}\color{white}\bfseries Criterio & \color{white}\bfseries Lo logré & \color{white}\bfseries En proceso & \color{white}\bfseries Necesito apoyo\\[.6mm]
Comprensión del tema & \milcopcion{Explico las ideas con mis palabras.} & \milcopcion{Reconozco las ideas, falta precisión.} & \milcopcion{Necesito repasar lo básico.}\\[.8mm]
\rowcolor{milcGris}Pensamiento computacional & \milcopcion{Usé los pilares al armar mi producto.} & \milcopcion{Usé algunos pilares.} & \milcopcion{Necesito releer la estación 2.}\\[.8mm]
Aplicación y producto & \milcopcion{Mi producto cumple los criterios.} & \milcopcion{Está armado, con ajustes pendientes.} & \milcopcion{Aún está en construcción.}\\[.8mm]
\rowcolor{milcGris}Reflexión 5 dimensiones & \milcopcion{Respondí cada una con honestidad.} & \milcopcion{Respondí algunas con detalle.} & \milcopcion{Mis respuestas son muy generales.}\\[.8mm]
Triángulo de pensamiento & \milcopcion{Conecté las 3 voces con mi trabajo.} & \milcopcion{Conecté al menos una.} & \milcopcion{No las relaciono con mi proceso.}\\[.8mm]
\end{tabularx}}

\vspace{3mm}
\textbf{Hoy entendí que:}\\[1mm]
\linead{\linewidth}\\[1mm]
\linead{\linewidth}

\textbf{Mi compromiso para la próxima sesión es:}\\[1mm]
\linead{\linewidth}\\[1mm]
\linead{\linewidth}

\begin{softbox}{milcMostaza}{milcGris}{Cita de despedida · Marco Aurelio}
\textit{``El obstáculo en el camino se vuelve el camino. Lo que estorba al camino se vuelve camino.''}\\[1mm]
Cada error de hoy, bien observado, se convierte en pieza del camino que pisarás mañana.
\end{softbox}

\small
\begin{tabularx}{\textwidth}{>{\bfseries}p{3.6cm}Y}
\rowcolor{milcGradeDark!12}Nombre completo: & \linead{10cm}\\
Fecha: & \linead{4cm} \quad Curso/grupo: \linead{4cm}\\
Mi firma: & \linead{9cm}\\
\end{tabularx}
\normalsize

\begin{infoband}{milcGradeDark}
\textbf{Para la próxima sesión.} Llega con el capítulo de muestra diagramado y el archivo de plantilla aplicado. En la sesión 9 vas a integrar \textbf{todos los capítulos en el PDF final de 80 páginas} con portada incluida. La diagramación de hoy es el patrón; la producción final de mañana lo aplica al libro completo.
\end{infoband}

\end{document}
